% © 2026 Ing. David Jaroš — CC BY-NC-ND 4.0
%
% This work is licensed under a Creative Commons Attribution-NonCommercial-NoDerivatives
% 4.0 International License (CC BY-NC-ND 4.0).
%
% License History: Earlier drafts (up to v0.3) were released under CC BY 4.0.
% From v0.4 onward, all material is released under CC BY-NC-ND 4.0 to protect
% the integrity of the theoretical work during ongoing academic development.
%
% See LICENSE.md for full license text.

% bootstrap_step_m2_4point.tex
% Modular Bootstrap Step M2: 4-point function computation
%
% Goal: Compute ⟨Θ(z₁)Θ(z₂)Θ(z₃)Θ(z₄)⟩ on T² for the lowest charged mode.
% Status: IN PROGRESS — see MODULAR_BOOTSTRAP_K1_PLAN.md §4 Step M2
% Prerequisite: bootstrap_step_m1_conformal.tex (c=3, operator content)

\ifdefined\INCLUDEMODE
\else
  \documentclass[12pt,a4paper]{article}
  \usepackage[a4paper, margin=2.5cm]{geometry}
  \usepackage{amsmath, amssymb, amsthm}
  \usepackage{mathtools}
  \usepackage{hyperref}
  \usepackage{xcolor}
  \usepackage{mdframed}
  \usepackage{booktabs}
  \usepackage{tcolorbox}

  \newtheorem{theorem}{Theorem}[section]
  \newtheorem{lemma}[theorem]{Lemma}
  \newtheorem{proposition}[theorem]{Proposition}
  \newtheorem{corollary}[theorem]{Corollary}
  \theoremstyle{definition}
  \newtheorem{definition}[theorem]{Definition}
  \theoremstyle{remark}
  \newtheorem{remark}[theorem]{Remark}

  \newmdenv[backgroundcolor=yellow!20, linecolor=orange, linewidth=1.5pt]{gapbox}
  \newmdenv[backgroundcolor=green!15, linecolor=green!60!black, linewidth=1.5pt]{resultbox}
  \newmdenv[backgroundcolor=red!8,   linecolor=red!60,   linewidth=1.5pt]{warnbox}
  \newmdenv[backgroundcolor=blue!8,  linecolor=blue!50,  linewidth=1.5pt]{insightbox}

  \title{\textbf{Modular Bootstrap Step M2}\\
         \large 4-Point Function of the UBT Biquaternion Field on $T^2$}
  \author{Ing.\ David Jaro\v{s}}
  \date{2026-04-28}

  \begin{document}
  \maketitle
\fi

%──────────────────────────────────────────────────────────────────────────────
\section{Purpose and Status}
\label{sec:m2:purpose}
%──────────────────────────────────────────────────────────────────────────────

\begin{tcolorbox}[colback=blue!4!white, colframe=blue!50!black,
                  title=Step M2 — 4-Point Function Computation]
\begin{tabular}{ll}
Parent plan      & \texttt{MODULAR\_BOOTSTRAP\_K1\_PLAN.md} \\
Prerequisite     & \texttt{bootstrap\_step\_m1\_conformal.tex}: $c = 3$, vertex operators \\
Goal             & Explicit $\langle V V V V \rangle$ as function of cross-ratio $\eta$ \\
Milestone        & 2026-05-12 \\
Status           & IN PROGRESS \\
Acceptance       & Closed-form 4-point function; charge conservation enforced \\
\end{tabular}
\end{tcolorbox}

\begin{warnbox}
\textbf{No-fit rule}: No constant fitted to $\alpha$ or $m_e$.  All inputs
from M1 (UBT axioms + free-boson CFT at self-dual radius).
\end{warnbox}

%──────────────────────────────────────────────────────────────────────────────
\section{Vertex Operators for the Charged Sector}
\label{sec:m2:vertex}
%──────────────────────────────────────────────────────────────────────────────

From Step M1 (\texttt{bootstrap\_step\_m1\_conformal.tex}), the CFT is three
free compact bosons $X^i$ ($i = 1,2,3$) at the self-dual radius
$R = R_\psi = \sqrt{\alpha'}$.  The vertex operators for KK/winding modes are:
\begin{equation}
  V_{m,n}(z, \bar{z})
  = \exp\!\Bigl[i\,p_L\, X_L(z) + i\,p_R\, X_R(\bar{z})\Bigr],
  \label{eq:m2:vertex_op}
\end{equation}
where (at the self-dual radius, per boson component):
\begin{equation}
  p_L = m + n, \quad p_R = m - n, \qquad m, n \in \mathbb{Z}.
  \label{eq:m2:momenta}
\end{equation}
The conformal dimensions are:
\begin{equation}
  h = \tfrac{1}{4}(m+n)^2, \qquad \bar{h} = \tfrac{1}{4}(m-n)^2.
  \label{eq:m2:dimensions}
\end{equation}

\paragraph{UBT constraint on the winding sector.}
From the Dirac quantisation condition on $S^1_\psi$
(canonical/appendices/appendix\_alpha\_geometry.tex \S1), the winding number
$n$ on the ψ-circle is the \emph{electromagnetic winding number} of the
Θ field.  The UBT gauge field $A_\psi$ has minimal charge $e$, so the
physically relevant charged vertex operators are $V_{m,n}$ with $|n| = 1$
(unit winding, minimal charge).  We take the \textbf{lowest charged mode}:
\begin{equation}
  V \equiv V_{0,1}(z, \bar{z}),
  \quad h = \bar{h} = \tfrac{1}{4},
  \quad \text{spin } s = h - \bar{h} = 0.
  \label{eq:m2:lowestmode}
\end{equation}
This operator has dimension $\Delta = h + \bar{h} = 1/2$, which coincides with
the dimension of the $j = 1/2$ primary of $SU(2)_1$ WZW (the spinor primary).
This is the key operator for the crossing symmetry test in Step M3.

%──────────────────────────────────────────────────────────────────────────────
\section{4-Point Function: Free-Boson Formula}
\label{sec:m2:4point}
%──────────────────────────────────────────────────────────────────────────────

The 4-point function on the sphere $\mathbb{CP}^1$ for four unit-winding
vertex operators $V_{0,1}$ of a single free compact boson is, using the
Coulomb-gas representation and charge conservation
($\sum_i n_i = 0$; place two operators with $n = +1$ and two with $n = -1$):
\begin{align}
  \bigl\langle V_{0,1}(z_1) V_{0,1}(z_2) V_{0,-1}(z_3) V_{0,-1}(z_4)
  \bigr\rangle
  &= \prod_{i < j} |z_i - z_j|^{2 p_{L,i} p_{L,j}}
  \label{eq:m2:4point_raw}
\end{align}
with $p_{L,1} = p_{L,2} = +1$ and $p_{L,3} = p_{L,4} = -1$ (and similarly for
$p_R$).  This gives:
\begin{align}
  \bigl\langle V_{0,1}(z_1) V_{0,1}(z_2) V_{0,-1}(z_3) V_{0,-1}(z_4)
  \bigr\rangle
  &= \frac{|z_{12}|^{2} |z_{34}|^{2}}{|z_{13}|^{2} |z_{14}|^{2}
                                         |z_{23}|^{2} |z_{24}|^{2}}
    \cdot |z_{12}|^0 |z_{34}|^0,
  \label{eq:m2:4point_boson}
\end{align}
where $z_{ij} = z_i - z_j$.

\paragraph{Standard positions.}
Fix $z_1 = \infty$, $z_2 = 1$, $z_4 = 0$, $z_3 = z$ (standard gauge).
The cross-ratio is $\eta = z_{12}z_{34}/(z_{13}z_{24})$.  In this gauge:
\begin{equation}
  \bigl\langle V_{0,1}(\infty) V_{0,1}(1) V_{0,-1}(z) V_{0,-1}(0)
  \bigr\rangle
  = |z|^{-2h_{3}} |1-z|^{2 p_{L,2} p_{L,3}}
  = |z|^{-1/2} |1-z|^{-1/2}.
  \label{eq:m2:4point_gauge}
\end{equation}
Here the exponents $-1/2 = 2 p_{L,i} p_{L,j}$ with $|p_L| = 1$; the
$|z|^{-1/2}$ factor is the standard free-boson result for dimension-$1/4$
operators.

\paragraph{Three-component system.}
For three decoupled bosons (the full $\operatorname{Im}(\mathbb{H})$ sector),
each contributes independently.  The 4-point function of a charged mode
$\Theta^{(i)}$ is the product:
\begin{equation}
  \bigl\langle \Theta(\infty)\, \Theta(1)\, \Theta(z)\, \Theta(0)
  \bigr\rangle_{c=3}
  = \Bigl(|z|^{-1/2} |1-z|^{-1/2}\Bigr)^{\!3}
  = |z|^{-3/2} |1-z|^{-3/2}.
  \label{eq:m2:4point_c3}
\end{equation}

%──────────────────────────────────────────────────────────────────────────────
\section{OPE Decomposition}
\label{sec:m2:ope}
%──────────────────────────────────────────────────────────────────────────────

The 4-point function~\eqref{eq:m2:4point_c3} admits an OPE decomposition in the
$s$-channel ($z \to 0$, operators at $z$ and $0$ fused):
\begin{equation}
  \bigl\langle \Theta \Theta \Theta \Theta \bigr\rangle
  = \sum_{\mathcal{O}} C_{\Theta\Theta\mathcal{O}}^2
    \, \mathcal{F}_{\mathcal{O}}(z) \, \overline{\mathcal{F}}_{\mathcal{O}}(\bar{z}),
  \label{eq:m2:ope_sum}
\end{equation}
where $\mathcal{F}_{\mathcal{O}}$ is the Virasoro conformal block for exchange
of operator $\mathcal{O}$ with dimension $\Delta_{\mathcal{O}} = h_{\mathcal{O}} + \bar{h}_{\mathcal{O}}$.

For the free-boson result~\eqref{eq:m2:4point_c3}, the $s$-channel OPE
($z \to 0$) gives the operator $\mathcal{O}_0$ with $h = 0$ (identity) from the
$n = 0$ sector and $\mathcal{O}_{1/2}$ with $h = 1/4$ from the $n = \pm 1$
winding sector.  The OPE coefficient is:
\begin{equation}
  C_{\Theta\Theta\mathcal{O}_{1/2}}
  = 1 \quad \text{(free-boson normalisation)}.
  \label{eq:m2:ope_coeff}
\end{equation}
This result is passed to Step M3 as the \emph{input} for crossing symmetry.

%──────────────────────────────────────────────────────────────────────────────
\section{Step M2 Result and Handoff to Step M3}
\label{sec:m2:result}
%──────────────────────────────────────────────────────────────────────────────

\begin{resultbox}
\textbf{Step M2 result}:
\begin{equation}
  \bigl\langle \Theta(\infty)\, \Theta(1)\, \Theta(z)\, \Theta(0)
  \bigr\rangle_{c=3}
  = |z|^{-3/2}\, |1 - z|^{-3/2}.
  \label{eq:m2:final}
\end{equation}
The OPE intermediate states are the identity ($h = 0$) and the unit-winding
operator ($h = \bar{h} = 1/4$, $\Delta = 1/2$).

\textbf{Circularity check}: The result uses only
$\dim_{\mathbb{R}}(\operatorname{Im}\mathbb{H}) = 3$ and the minimal-charge
winding sector $|n| = 1$.  No $\alpha$, no $m_e$.  [CLEAN]

\textbf{Key question for Step M3}: Does the function $|z|^{-3/2}|1-z|^{-3/2}$
satisfy crossing symmetry $z \leftrightarrow 1-z$ with the same OPE coefficient?
If yes: which values of $k$ are consistent?
\end{resultbox}

\begin{gapbox}
\textbf{Open question}: The 4-point function~\eqref{eq:m2:final} is manifestly
symmetric under $z \leftrightarrow 1-z$ (since
$|(z)(1-z)|^{-3/2} = |(1-z)(z)|^{-3/2}$), so crossing symmetry is automatic
for the full free-boson theory.  The \emph{non-trivial} constraint arises
when asking whether the UBT operator spectrum is \emph{truncated} (Step M4:
only $j = 0, 1/2$ primaries from $SU(2)_1$, vs.\ the full winding tower).
Step M3 makes this precise.
\end{gapbox}

\ifdefined\INCLUDEMODE
\else
  \end{document}
\fi
